\documentclass[12pt,letterpaper]{article}
\usepackage{graphicx,textcomp}
\usepackage{natbib}
\usepackage{setspace}
\usepackage{fullpage}
\usepackage{color}
\usepackage[reqno]{amsmath}
\usepackage{amsthm}
\usepackage{fancyvrb}
\usepackage{amssymb,enumerate}
\usepackage[all]{xy}
\usepackage{endnotes}
\usepackage{lscape}
\newtheorem{com}{Comment}
\usepackage{float}
\usepackage{hyperref}
\newtheorem{lem} {Lemma}
\newtheorem{prop}{Proposition}
\newtheorem{thm}{Theorem}
\newtheorem{defn}{Definition}
\newtheorem{cor}{Corollary}
\newtheorem{obs}{Observation}
\usepackage[compact]{titlesec}
\usepackage{dcolumn}
\usepackage{tikz}
\usetikzlibrary{arrows}
\usepackage{multirow}
\usepackage{xcolor}
\newcolumntype{.}{D{.}{.}{-1}}
\newcolumntype{d}[1]{D{.}{.}{#1}}
\definecolor{light-gray}{gray}{0.65}
\usepackage{url}
\usepackage{listings}
\usepackage{color}

\definecolor{codegreen}{rgb}{0,0.6,0}
\definecolor{codegray}{rgb}{0.5,0.5,0.5}
\definecolor{codepurple}{rgb}{0.58,0,0.82}
\definecolor{backcolour}{rgb}{0.95,0.95,0.92}

\lstdefinestyle{mystyle}{
	backgroundcolor=\color{backcolour},   
	commentstyle=\color{codegreen},
	keywordstyle=\color{magenta},
	numberstyle=\tiny\color{codegray},
	stringstyle=\color{codepurple},
	basicstyle=\footnotesize,
	breakatwhitespace=false,         
	breaklines=true,                 
	captionpos=b,                    
	keepspaces=true,                 
	numbers=left,                    
	numbersep=5pt,                  
	showspaces=false,                
	showstringspaces=false,
	showtabs=false,                  
	tabsize=2
}
\lstset{style=mystyle}
\newcommand{\Sref}[1]{Section~\ref{#1}}
\newtheorem{hyp}{Hypothesis}


\title{Problem Set 4}
\date{Due: November 26, 2026}
\author{Ethan Kelly Zeh (26359218)}


\begin{document}
	\maketitle
	\section*{Instructions}
	\begin{itemize}
		\item Please show your work! You may lose points by simply writing in the answer. If the problem requires you to execute commands in \texttt{R}, please include the code you used to get your answers. Please also include the \texttt{.R} file that contains your code. If you are not sure if work needs to be shown for a particular problem, please ask.
		\item Your homework should be submitted electronically on GitHub.
		\item This problem set is due before 23:59 on Thursday November 26, 2026. No late assignments will be accepted.
	\end{itemize}



	\vspace{.5cm}
\section*{Question 1: Economics}
\vspace{.25cm}
\noindent 	
In this question, use the \texttt{prestige} dataset in the \texttt{car} library. First, run the following commands:

\begin{verbatim}
install.packages(car)
library(car)
data(Prestige)
help(Prestige)
\end{verbatim} 


\noindent We would like to study whether individuals with higher levels of income have more prestigious jobs. Moreover, we would like to study whether professionals have more prestigious jobs than blue and white collar workers.

\newpage
\begin{enumerate}
	
	\item [(a)]
	Create a new variable \texttt{professional} by recoding the variable \texttt{type} so that professionals are coded as $1$, and blue and white collar workers are coded as $0$ (Hint: \texttt{ifelse}).
	\vspace{.5cm}

	\textbf{Explanation}
	\vspace{.5cm}

	After loading the \texttt{Prestige} data, I use \texttt{ifelse} to code \texttt{prof} as 1 and \texttt{bc}/\texttt{wc} as 0. The table checks the recode: 31 professional occupations are coded 1 and the 67 blue and white collar occupations are coded 0. The 4 occupations with no \texttt{type} stay \texttt{NA}, so \texttt{lm} drops them later and the models below use 98 observations.
	\vspace{1cm}

	\textbf{Code}
	\lstinputlisting[language=R, firstline=33, lastline=41]{PS04_answersEZ.R}
	\vspace{1cm}
	\textbf{Output}
	\begin{verbatim}
        0  1 <NA>
  bc   44  0    0
  prof  0 31    0
  wc   23  0    0
  <NA>  0  0    4
	\end{verbatim}

	\vspace{6cm}
	
	
	\item [(b)]
	Run a linear model with \texttt{prestige} as an outcome and \texttt{income}, \texttt{professional}, and the interaction of the two as predictors (Note: this is a continuous $\times$ dummy interaction.)
	\vspace{.5cm}

	\textbf{Explanation}
	\vspace{.5cm}

	I fit the model below 
	\[
	\texttt{prestige}_i = \beta_0 + \beta_1\,\texttt{income}_i + \beta_2\,\texttt{professional}_i + \beta_3\,(\texttt{income}_i \times \texttt{professional}_i) + \varepsilon_i
	\]
	with \texttt{lm}, using \texttt{income * professional} which adds both main effects and the interaction effect. All four coefficients are significant at the 0.05 level. The interaction is negative ($-0.0023$, $p < 0.001$), so income appears to matter less for prestige among professionals than for non.
	\vspace{1cm}

	\textbf{Code}
	\lstinputlisting[language=R, firstline=42, lastline=50]{PS04_answersEZ.R}
	\vspace{1cm}

	\textbf{Stargazer table/Output}
	% Table created by stargazer v.5.2.3 by Marek Hlavac, Social Policy Institute. E-mail: marek.hlavac at gmail.com
\begin{table}[H] \centering 
  \caption{Prestige regressed on income, professional and their interaction} 
  \label{} 
\begin{tabular}{@{\extracolsep{5pt}}lc} 
\\[-1.8ex]\hline 
\hline \\[-1.8ex] 
 & \multicolumn{1}{c}{\textit{Dependent variable:}} \\ 
\cline{2-2} 
\\[-1.8ex] & Prestige \\ 
\hline \\[-1.8ex] 
 Income & 0.0032$^{***}$ \\ 
  & (0.0005) \\ 
  & \\ 
 Professional & 37.7813$^{***}$ \\ 
  & (4.2483) \\ 
  & \\ 
 Income $\times$ Professional & $-$0.0023$^{***}$ \\ 
  & (0.0006) \\ 
  & \\ 
 Constant & 21.1423$^{***}$ \\ 
  & (2.8044) \\ 
  & \\ 
\hline \\[-1.8ex] 
Observations & 98 \\ 
R$^{2}$ & 0.7872 \\ 
Adjusted R$^{2}$ & 0.7804 \\ 
\hline 
\hline \\[-1.8ex] 
\textit{Note:}  & \multicolumn{1}{r}{$^{*}$p$<$0.1; $^{**}$p$<$0.05; $^{***}$p$<$0.01} \\ 
\end{tabular} 
\end{table} 

	\vspace{6cm}
	\item [(c)]
	Write the prediction equation based on the result.
	\vspace{.5cm}

	Using the coefficients from above:
	\[
	\widehat{\texttt{prestige}}_i = 21.142 + 0.0032 \times \texttt{income}_i + 37.781 \times \texttt{professional}_i - 0.0023 \times (\texttt{income}_i \times \texttt{professional}_i)
	\]

	\vspace{.5cm}

	\textbf{Code}
	\lstinputlisting[language=R, firstline=51, lastline=51]{PS04_answersEZ.R}
	\vspace{.5cm}
	\textbf{Output}
	\begin{verbatim}
        (Intercept)              income        professional
        (Intercept)              income        professional income:professional 
       21.142258854         0.003170909        37.781279955        -0.002325709 
	\end{verbatim}
	
\newpage
	\item [(d)]
	Interpret the coefficient for \texttt{income}.
	\vspace{.5cm}

	The coefficient on \texttt{income} ($0.0032$) is the effect of income only when \texttt{professional} $= 0$, so for blue and white collar occupations (because of the interaction). Among non-professionals, each unit/\$1 increase in average income is associated with an increase of approx. 0.0032 points in the prestige score, which is approx. 3.2 points for every \$1,000. It is significant ($p < 0.001$).

	\vspace{5cm}	
	\item [(e)]
	Interpret the coefficient for \texttt{professional}.
	\vspace{.5cm}

	The coefficient on \texttt{professional} ($37.78$) is the difference in predicted presetige between professional and non-professional occupations only when \texttt{income} $= 0$. The interpretation of this is: at zero income, professional jobs are predicted to have a prestige score approx. 37.78 points higher. This is not particularly useful because no occupation has an average income of \$0. The gap between the two groups gets smaller as income increases as the interaction is negative.
	
	\newpage
	\item [(f)]
	What is the effect of a \$1,000 increase in income on prestige score for professional occupations? In other words, we are interested in the marginal effect of income when the variable \texttt{professional} takes the value of $1$. Calculate the change in $\hat{y}$ associated with a \$1,000 increase in income based on your answer for (c).
	\vspace{.5cm}

	\textbf{Explanation}
	\vspace{.5cm}

	When \texttt{professional} $= 1$, the slope on income is $\beta_1 + \beta_3$:
	\[
	\frac{\partial \widehat{\texttt{prestige}}}{\partial \texttt{income}} = \beta_1 + \beta_3 \times 1 = 0.0031709 - 0.0023257 = 0.0008452
	\]
	for a \$1,000 increase:
	\[
	\Delta \hat{y} = 1000 \times 0.0008452 \approx 0.845
	\]
	Therefore, in regards to professional occupations, a \$1,000 increase in income is associated with an increase of approx. 0.85 points in the 'prestige' score, much smaller than the approx. 3.2 points for non-professionals.
	\vspace{1cm}

	\textbf{Code}
	\lstinputlisting[language=R, firstline=54, lastline=58]{PS04_answersEZ.R}
	\vspace{.5cm}
	\textbf{Output}
	\begin{verbatim}
   income
0.8452004
	\end{verbatim}

	\vspace{10cm}
	
	
	\item [(g)]
	What is the effect of changing one's occupations from non-professional to professional when her income is \$6,000? We are interested in the marginal effect of professional jobs when the variable \texttt{income} takes the value of $6,000$. Calculate the change in $\hat{y}$ based on your answer for (c).
	\vspace{.5cm}

	\textbf{Explanation}
	\vspace{.5cm}

	When \texttt{professional} switches from 0 to 1, $\hat{y}$ changes by $\beta_2 + \beta_3 \times \texttt{income}$. With \texttt{income} $= 6000$:
	\[
	\Delta \hat{y} = \beta_2 + \beta_3 \times 6000 = 37.7813 + (-0.0023257 \times 6000) = 37.7813 - 13.9542 \approx 23.83
	\]
	So at an income of \$6,000, moving from a non-professional to a professional occupation is associated with an increase of approx. 23.8 points in the prestige score. This is smaller than the 37.8 point gap at zero income because the negative interaction shrinks the professional premium as income rises.
	\vspace{1cm}

	\textbf{Code}
	\lstinputlisting[language=R, firstline=59, lastline=62]{PS04_answersEZ.R}
	\vspace{.5cm}
	\textbf{Output}
	\begin{verbatim}
professional
    23.82708
	\end{verbatim}


	
\end{enumerate}

\newpage

\section*{Question 2: Political Science}
\vspace{.25cm}
\noindent 	Researchers are interested in learning the effect of all of those yard signs on voting preferences.\footnote{Donald P. Green, Jonathan	S. Krasno, Alexander Coppock, Benjamin D. Farrer,	Brandon Lenoir, Joshua N. Zingher. 2016. ``The effects of lawn signs on vote outcomes: Results from four randomized field experiments.'' Electoral Studies 41: 143-150. } Working with a campaign in Fairfax County, Virginia, 131 precincts were randomly divided into a treatment and control group. In 30 precincts, signs were posted around the precinct that read, ``For Sale: Terry McAuliffe. Don't Sellout Virgina on November 5.'' \\

Below is the result of a regression with two variables and a constant.  The dependent variable is the proportion of the vote that went to McAuliff's opponent Ken Cuccinelli. The first variable indicates whether a precinct was randomly assigned to have the sign against McAuliffe posted. The second variable indicates
a precinct that was adjacent to a precinct in the treatment group (since people in those precincts might be exposed to the signs).  \\

\vspace{.5cm}
\begin{table}[!htbp]
	\centering 
	\textbf{Impact of lawn signs on vote share}\\
	\begin{tabular}{@{\extracolsep{5pt}}lccc} 
		\\[-1.8ex] 
		\hline \\[-1.8ex]
		Precinct assigned lawn signs  (n=30)  & 0.042\\
		& (0.016) \\
		Precinct adjacent to lawn signs (n=76) & 0.042 \\
		&  (0.013) \\
		Constant  & 0.302\\
		& (0.011)
		\\
		\hline \\
	\end{tabular}\\
	\footnotesize{\textit{Notes:} $R^2$=0.094, N=131}
\end{table}

\vspace{.5cm}
\begin{enumerate}
	\item [(a)] Use the results from a linear regression to determine whether having these yard signs in a precinct affects vote share (e.g., conduct a hypothesis test with $\alpha = .05$).
	\vspace{.5cm}

	\textbf{Explanation}
	\vspace{.5cm}

	The following are the two hypothesis:

	$H_0: \beta = 0$ (the yard signs have no effect on cuccinelli's vote share) \\
	$H_A: \beta \neq 0$ \\

	With $N = 131$ and 2 predictors: $df = 131 - 2 - 1 = 128$. We then calculate the test statistic:
	\[
	t = \frac{\hat{\beta}}{SE(\hat{\beta})} = \frac{0.042}{0.016} = 2.63
	\]
	The two-tailed critical value is $t_{0.975,\,128} = 1.979$ and $p \approx 0.0097$. Since $2.625 > 1.979$ ($p < 0.05$), we reject $H_0$. Therefore, precincts assigned the lawn signs appear to have a Cuccinelli vote share approx. 4.2 percentage points higher than the control precincts (those without the signs). As this was a valid, randomly assigned expirement, this can be read as a causal effect.
	\vspace{1cm}

	\textbf{Code}
	\lstinputlisting[language=R, firstline=63, lastline=71]{PS04_answersEZ.R}
	\vspace{.5cm}
	\textbf{Output}
	\begin{verbatim}
t = 2.625 | p = 0.0097 | critical t = 1.979
	\end{verbatim}

	\newpage
	\item [(b)]  Use the results to determine whether being
	next to precincts with these yard signs affects vote
	share (e.g., conduct a hypothesis test with $\alpha = .05$).
	\vspace{.5cm}

	\textbf{Explanation}
	\vspace{.5cm}

	$H_0: \beta_{\text{adjacent}} = 0$ (being next to a precinct with signs has no effect on vote share) \\
	$H_A: \beta_{\text{adjacent}} \neq 0$ \\

	Again with $df = 128$:
	\[
	t = \frac{\hat{\beta}_{\text{adjacent}}}{SE(\hat{\beta}_{\text{adjacent}})} = \frac{0.042}{0.013} \approx 3.231
	\]
	Since $3.23 > 1.979$ ($p \approx 0.0016 < 0.05$), we reject fhe null $H_0$. Precincts next to the treated precincts also appear to have a Cuccinelli vote share approx. 4.2 percentage points higher than the control precincts. This suggests the signs have a sort-of spill-over effect on voters in neighbouring precincts who may also see them. 	
	
	\vspace{1cm}

	\textbf{Code}
	\lstinputlisting[language=R, firstline=72, lastline=77]{PS04_answersEZ.R}
	\vspace{.5cm}
	\textbf{Output}
	\begin{verbatim}
t = 3.231 | p = 0.0016 | critical t = 1.979
	\end{verbatim}

	\vspace{7cm}
	\item [(c)] Interpret the coefficient for the constant term substantively.
	\vspace{.5cm}

	The constant ($0.30$) is the predicted proportion of the vote for Cuccinelli when both variables are 0, which means precincts that were neither assigned signs nor next to a precinct with signs. In these precincts Cuccinelli is predicted to get approx. 30.2\% of the vote. This is the baseline the effects in (a) and (b) are compared against. It is significant ($t = 0.302/0.011 \approx 27.5$), but that only means the baseline vote share is not 0.
	\vspace{2cm}
	
	\item [(d)] Evaluate the model fit for this regression.  What does this	tell us about the importance of yard signs versus other factors that are not modeled?
	\vspace{.5cm}

	The $R^2$ is $0.094$, so the sign variables explain only about 9.4\% of the variation in Cuccinelli's vote share. So the yard signs appear to have a real and statistically significant effect of about 4 percentage points, but they are a small part of what determines vote share compared to other factors. A low $R^2$ does not make the estimates in (a) and (b) wrong: because the signs were randomly assigned, the omitted factors should not be correlated with treatment, so the coefficients are still unbiased. The model just does not predict vote share that well overall; it has another usecase.

\end{enumerate}  


\end{document}
